\documentclass[11pt]{report}
\usepackage[utf8]{inputenc}
\usepackage{amsmath, amssymb}
\usepackage{listings}
\usepackage{xcolor}
\usepackage{a4wide}
\usepackage{epsfig}
\usepackage{amssymb}
\usepackage{amsmath}
\usepackage{minted}
\usepackage{enumerate}
\usepackage{fancyvrb}
\usepackage{alltt}
\usepackage{fleqn}
\usepackage{epic}
\usepackage{color} 
\usepackage{theorem}
\usepackage{hyperref}
\usepackage{tikz}
\usepackage[all]{hypcap}
\hypersetup{
        colorlinks = true, % comment this to make xdvi work
        linkcolor  = blue,
        citecolor  = red,
        filecolor  = [rgb]{0.1, 0.1, 1.0},
        urlcolor   = [rgb]{0.7, 0.0, 0.7},
        pdfborder  = {0 0 0} 
}

\usepackage{fancyhdr}
\usepackage{lastpage} 

\definecolor{bg}{rgb}{0.9,0.8,0.6}
\definecolor{amber}{rgb}{1.0, 0.75, 0.0}
\definecolor{gold}{rgb}{1.0, 0.843, 0.0}

\renewcommand*{\familydefault}{\sfdefault}

% Set up listings for Python
\lstset{
    language=Python,
    basicstyle=\ttfamily\small,
    keywordstyle=\color{blue},
    stringstyle=\color{red},
    commentstyle=\color{green!60!black},
    showstringspaces=false,
    breaklines=true,
    frame=single,
    captionpos=b
}



\pagestyle{fancy}

\fancyfoot[C]{--- \thepage/\pageref{LastPage}\ ---}

\fancypagestyle{plain}{%
\fancyhf{}
\fancyfoot[C]{--- \thepage/\pageref{LastPage}\ ---}
\renewcommand{\headrulewidth}{0pt}
}

\renewcommand{\chaptermark}[1]{\markboth{\chaptername \ \thechapter.\ #1}{}}
\renewcommand{\sectionmark}[1]{\markright{\thesection. \ #1}{}}
\fancyhead[R]{\leftmark}
\fancyhead[L]{\rightmark}

\definecolor{amethyst}{rgb}{0.2, 0.4, 0.6}
\definecolor{orange}{rgb}{1, 0.9, 0.0}

{\theorembodyfont{\sf}
\newtheorem{Definition}{Definition}
\newtheorem{Axiom}[Definition]{Axiom}
\newtheorem{Notation}[Definition]{Notation}
\newtheorem{Theorem}[Definition]{Theorem}
}

\newcommand{\blue}[1]{{\color{blue}#1}}


\newcommand{\qed}{\hspace*{\fill} $\Box$
\vspace*{0.2cm}

}

\newcommand{\eod}{\hspace*{\fill} $\diamond$}
\newcommand{\eox}{\hspace*{\fill} $\diamond$}

\newcommand{\dangersign}[1][Danger]{
    \begin{tikzpicture}[scale=0.7]
        % Draw the triangle
        \draw[thick, red, fill=yellow] (0,0) -- (2,3.46) -- (4,0) -- cycle;
        % Add the text (default is "Danger")
        \node at (2,0.8) {\textbf{\textsc{#1}}};
    \end{tikzpicture}
}

\newcommand{\wip}[1][\ \ Work in Progress!]{
    \begin{tikzpicture}[scale=0.7]
        % Draw rectancle
        \draw[thick, blue, fill=orange] (0,3) -- (6.5, 3) -- (6.5,4) -- (0, 4) -- cycle;
        \node at (3,3.5) {\textbf{\textsc{#1}}};
    \end{tikzpicture}
}
\title{\epsfig{file=dhbw-logo.pdf, scale=1.5}\\[0.3cm]  
       Functional Programming in \textsl{Haskell} \\[0.3cm]
      --- Winter 2025 ---}
\author{Prof.~Dr.~Karl Stroetmann}
\date{\today \\[1.5cm]
  \dangersign \wip \dangersign
}

\newcommand{\ds}{\displaystyle}
\newcommand{\solution}{\vspace*{0.2cm}

\noindent
\textbf{Solution}: }

\newcounter{aufgabe}
\newcommand{\exercise}{\vspace*{0.2cm}
\stepcounter{aufgabe}

\noindent
\textbf{Exercise \arabic{aufgabe}}: }

\newcommand{\exerciseStar}{\vspace*{0.2cm}
\stepcounter{aufgabe}

\noindent
\textbf{Exercise \arabic{aufgabe}$^*$}: }


\newcommand{\example}{\vspace*{0.2cm}

\noindent
\textbf{Example}: \ }

\newcommand{\examples}{\vspace*{0.2cm}

\noindent
\textbf{Examples}: \ }
 
\newcommand{\remark}{\vspace*{0.2cm}
\noindent
\textbf{Remark}: }

\newcommand{\lb}{\hspace*{\fill} \linebreak}


\def\pair(#1,#2){\langle #1, #2 \rangle}

\newlength{\mylength}
\setlength{\mathindent}{1.3cm}

\begin{document}
\maketitle

\begin{abstract}
\noindent
This is a very short introduction to functional programming with
\href{https://en.wikipedia.org/wiki/Haskell}{\textsl{Haskell}}. 
It is not intended to be a Haskell course.  My intention merely is for the reader to get a taste of what
programming in \textsl{Haskell} feels like.  If I succeed in convincing the reader that \textsl{Haskell} is a
programming language that is both usefull and mind extending, then I consider my job done.
Therefore, this short paper will just present some of the highlights of the programming language
\textsl{Haskell} via examples that I hope the reader finds intriguing enough so that she feels inclined to read
some of the outstanding books introducing Haskell in more depth.  There are three books that I recommend for
those readers who want to understand \textsl{Haskell} in more depth:
\begin{enumerate}
\item \emph{Effective Haskell: Solving Real-World Problems with Strongly Typed Functional Programming} 

      by Rebecca Skinner \cite{skinner:2023}.
\item \emph{Programming in Haskell}, second edition 

      by Graham Hutton \cite{hutton:2016}.
  
\item \emph{Learn You a Haskell for Great Good!},

      by Miran Lipova\v{c}a \cite{lipovaca:2011}.
\end{enumerate}
\end{abstract}

\chapter{Introduction}
In this introduction I will do two things:
\begin{enumerate}
\item First, I discuss those features of Haskell that set Haskell apart from other programming languages.
\item Second, I will present a few short example programs that give a first taste of Haskell.
\end{enumerate}

\section{Why Haskell is Different}
Before we present any details of \textsl{Haskell}, let us categorize this programming language so that we have
an idea about what to expect.  \textsl{Haskell} has the following properties:
\begin{enumerate}
\item \textsl{Haskell} is a \blue{functional programming language}.

  A functional programming language is any programming language that treats functions as
  \blue{first class citiziens}:
  \begin{enumerate}
  \item A function can be given as an argument to another function.
  \item A function can be produce a function as its result.
  \end{enumerate}
  A well known programming language that supports functional programming is \textsl{Python} and there are
  several books discussing functional programming in \textsl{Python}, e.g.~\cite{lott:2022}, \cite{mertz:2015},
  and \cite{reid:2023}.  

\item \textsl{Haskell} is \blue{statically typed}.

  Every variable in Haskell has a fixed type, which can not be changed.  In this respect, \textsl{Haskell} is
  similar to the programming language \textsl{Java}.  However, in contrast to \textsl{Java}, we do not have to
  declare the type of every variable and every function because most of the time the type of a variable can be
  \blue{inferred} by the type system.  Therefore, in Haskell we usually specify only the types of non-trivial
  functions.

  The benefit of this approach is that many type errors will already be caught by the compiler.  This is in
  contrast to programs written in a dynamically typed language like \textsl{Python}, where type errors are only
  discovered at runtime.

\item \textsl{Haskell} is a \blue{pure} functional programming language.

  Once a variable is assigned a value, this value can not be changed.  For example, if we want to sort a list,
  we are not able to change the list data structure.  All we can do is to compute a new list which contains the
  same elements as the old list and which, furthermore, is sorted.

  The property of being a \blue{pure} language sets \textsl{Haskell} apart from most other programming
  languages.  Even the language
  \href{https://en.wikipedia.org/wiki/Scala_(programming_language)}{\textsl{Scala}}, which is designed as a
  modern functional programming language, is not pure.

  What is the big deal about purity?  On one hand, it forces
  the user to program in a declarative style.  Although, in general, nobody likes to be forced to do something,
  there is a huge benefit in pure programming.
  \begin{enumerate}
  \item In a pure programming language, functions will always return the same result when they are called with
    the same arguments.  This property is called \blue{referential transparency}.
    This makes reasoning about code easier, as you can replace a function call with its result without changing
    the behavior of the program.  Therefore the correctness of functions can be verified mathematically.
  \item Since pure functions do not depend on or modify external state, their behavior is entirely predictable.
    The advantage is that testing becomes straightforward because functions can be tested in isolation without
    worrying about interactions with other parts of the system. 
  \item Compilers for pure languages can make aggressive optimizations, such as caching function results
    (\blue{memoization}) or reordering computations, because they know that functions are side-effect-free. 
    The advantage is that programs can often run faster.
  \item Furthermore, \blue{concurrency} becomes much easier to manage when functions do not use global variables and do
    not change their arguments.
  \end{enumerate}
  
\item \textsl{Haskell} is a \blue{compiled} language similar to \textsl{Java} and \texttt{C}, but additionally
  offers an interpreter.  Having an interpreter is beneficial for rapid prototyping.  The property that
  \textsl{Haskell} programs can be compiled ensures that the resulting programs can befaster than,
  for example, \textsl{Python} programs.
\item \textsl{Haskell} is a \blue{lazy} language.  The programming language \texttt{C} is an \blue{eager}
  language.  If an expression  of the form
  \\[0.2cm]
  \hspace*{1.3cm}
  $f(a_1, \cdots, a_n)$
  \\[0.2cm]
  has to be evaluated, first the subexpressions $a_1$, $\cdots$ $a_n$ are evaluated.  Let us assume that $a_i$
  is evaluated to to value $x_i$.  The, $f(x_1, \cdots, x_n)$ is computed.  This might be very inefficient.
  Consider the following contrived example:

\begin{minted}[ frame         = lines, 
                 framesep      = 0.3cm, 
                 firstnumber   = 1,
                 bgcolor       = bg,
                 numbers       = left,
                 numbersep     = -0.2cm,
                 xleftmargin   = 0.8cm,
                 xrightmargin  = 0.8cm,
               ]{C}
  int f(int x, int y) {
      if (x == 2) {
          return 42;
      }
      return 2 * y;
  }
\end{minted}
Let us assume that the expression
\\[0.2cm]
\hspace*{1.3cm}
\texttt{f(h(0), g(1))}
\\[0.2cm]
needs to be evaluated and that the computation of \texttt{g(1)} is very expensive.  In a
\texttt{C}-program, the expressions and \texttt{h(0)} and \texttt{g(1)} will be both evaluated.  If it turns out that
\texttt{h(0)} is \texttt{2}, then the evaluation of \texttt{g(1)} is not really necessary.  Nevertheless, in \texttt{C} this
evaluation takes place because \texttt{C} has an \blue{eager} evaluation strategy.  In contrast, an equivalent
\textsl{Haskell} would not evaluate the expression \texttt{h(0)} and hence would be much more efficient.

\item \textsl{Haskell} is difficult to learn.

  You might ask yourself why \textsl{Haskell} hasn't been adopted more widely.  After all, it has all these
  cool features mentioned above.  The reason is that learning \textsl{Haskell} is a lot more difficult then
  learning a language like \textsl{Python} or \textsl{Java}.  There are two reasons for this:
  \begin{enumerate}[(a)]
  \item First, \textsl{Haskell} differs a lot from those languages that most people know.
  \item In order to be very concise, the syntax of \textsl{Haskell} is quite different from the syntax of
        established programming languages.
  \item \textsl{Haskell} requires the programmer to think on a very high level of abstraction.
        Many students find this difficult.
  \item Lastly, and most importantly, \textsl{Haskell} supports the use of a number of concepts
        like, e.g.~\href{https://en.wikipedia.org/wiki/Functor}{functors} and
        \href{https://en.wikipedia.org/wiki/Monad_(category_theory)}{monads} from 
        \href{https://en.wikipedia.org/wiki/Category_theory}{category theory}. 
        It takes both time and mathematical maturity to really understand these concepts.

        If you really want to understand the depth of \textsl{Haskell}, you have to dive into those topics.
        That said, while you have to understands both functors and monads, you do not have to understand
        category theory.
  \item Fortunately, it is possible to become productive in \textsl{Haskell} without understanding 
        functors and monads.  Therefore, this lecture will focuss on those parts of \textsl{Haskell} that are more easily
        accessible. 
  \end{enumerate}
\end{enumerate}

\section{A First Taste of Haskell}
\subsubsection{Computing All Prime Numbers}
A \href{https://en.wikipedia.org/wiki/Prime_number}{prime number} is a natural number $p$ that is different from $1$
and that can not be written as a product of two natural numbers $a$ and $b$ that are both different from $1$.  If
we denote the set of all prime numbers with the symbol $\mathbb{P}$, we therefore have:
\\[0.2cm]
\hspace*{1.3cm}
$\ds \mathbb{P} = \bigl\{ p \in \mathbb{N} \mid n \not= 1 \wedge \forall a, b \in \mathbb{N}:(a \cdot b = p
\implies a = 1 \vee b = 1) \bigr\}$
\\[0.2cm]
An efficient method to compute the prime numbers is the
\href{https://en.wikipedia.org/wiki/Sieve_of_Eratosthenes}{sieve of Eratosthenes}.
This is an algorithm used to find all primes up to a given number. The method works by iteratively marking the
multiples of each prime number starting from 2. The numbers which remain unmarked at the end of the process are
the prime numbers.  For example, consider the list of integers from 2 to 30:
\\[0.2cm]
\hspace*{0.3cm}
$2, 3, 4, 5, 6, 7, 8, 9, 10, 11, 12, 13, 14, 15, 16, 17, 18, 19, 20, 21, 22, 23, 24, 25, 26, 27, 28, 29, 30$.
\\[0.2cm]
We will now compute the set of all primes less or equal than 30 using the Sieve of Eratosthenes.
\begin{enumerate}
\item The first number in this list is $2$ and hence $2$ is prime:
      \\[0.2cm]
      \hspace*{1.3cm}
      $\mathbb{P} = \{ 2, \cdots \}$.
\item We remove all multiples of 2 (i.e., 2, 4, 6, 8, \(\ldots\), 30).  This leaves us with the list
      \\[0.2cm]
      \hspace*{1.3cm}
      $3, 5, 7, 9, 11, 13, 15, 17, 19, 21, 23, 25, 27, 29$.
      \\[0.2cm]
      The first remaining number $3$ is prime.  Hence we have
      \\[0.2cm]
      \hspace*{1.3cm}
      $\mathbb{P} = \{ 2, 3, \cdots \}$
\item We remove all multiples of 3 (i.e., 3, 6, 9, 12, \(\ldots\), 27).  Note that some numbers (like 6 or 12,
      for instance) may already have been removed.  This leaves us with the list
      \\[0.2cm]
      \hspace*{1.3cm}
      $5, 7, 11, 13, 17, 19, 23, 25, 29$.
      \\[0.2cm]
      The first remaining number $5$ is prime. Hence we have
      \\[0.2cm]
      \hspace*{1.3cm}
      $\mathbb{P} = \{ 2, 3, 5, \cdots \}$.
\item We remove all multiples of 5 (i.e., 5, 10, 15, 20, \(\ldots\)). This leaves us with the list
      \\[0.2cm]
      \hspace*{1.3cm}
      $7, 11, 13, 17, 19, 23, 29$.
      \\[0.2cm]
      The first remaining number $7$ is prime. Hence we have
      \\[0.2cm]
      \hspace*{1.3cm}
      $\mathbb{P} = \{ 2, 3, 5, 7, \cdots \}$.
\item Since the square of $7$ is $49$, which is greater that 30, there is no need to remove the multiples of
      $7$,  since all multiples $a \cdot 7$ for $a < 7$ have already been removed and all multiples $a \cdot 7$
      for $a \geq 7$ are greater than 30. Hence the remaining numbers are all prime and we have found that the
      prime numbers less or equal than 30 are: 
      \\[0.2cm]
      \hspace*{1.3cm}
      $\{ 2, 3, 5, 7, 11, 13, 17, 19, 23, 29\}$. 
\end{enumerate}

\begin{figure}[!ht]
\centering
\begin{minted}[ frame         = lines, 
                 framesep      = 0.3cm, 
                 firstnumber   = 1,
                 bgcolor       = bg,
                 numbers       = left,
                 numbersep     = -0.2cm,
                 xleftmargin   = 0.8cm,
                 xrightmargin  = 0.8cm,
               ]{python3}
def sieve_of_eratosthenes(n):
    """Return a list of prime numbers up to n (inclusive)."""
    if n < 2:
        return []    
    is_prime = [True] * (n + 1)
    is_prime[0] = is_prime[1] = False  # 0 and 1 are not primes
    p = 2
    while p * p <= n:
        if is_prime[p]:
            for i in range(p * p, n + 1, p):
                is_prime[i] = False
        p += 1
    return [i for i, prime in enumerate(is_prime) if prime]
\end{minted}
\vspace*{-0.3cm}
\caption{A \textsl{Python program to compute the prime numbers up to $n$.}}
\label{fig:primes.py}
\end{figure}
Figure \ref{fig:primes.py} on page \pageref{fig:primes.py} shows a Python script that implements the Sieve of
Eratosthenes to compute all prime numbers up to a limit \texttt{n}.  This script first initializes the list
\texttt{is\_prime} to track prime status. It then iteratively marks the multiples of each prime number, finally
collecting and printing all numbers that remain marked as prime.  This script implements one optimization: If
\texttt{p} is a prime, then only the multiples of \texttt{p} that have the form $a \cdot \mathtt{p}$ with
$a \geq \mathtt{p}$ have to be removed from the list, since a product of the form $a \cdot \mathtt{p}$ with
$a < \mathtt{p}$ has already been removed when removing multiples  
of $a$ or, if $a$ is not prime, of multiples of whatever prime is contained in $a$.



\begin{figure}[!ht]
\centering
\begin{minted}[ frame         = lines, 
                 framesep      = 0.3cm, 
                 firstnumber   = 1,
                 bgcolor       = bg,
                 numbers       = left,
                 numbersep     = -0.2cm,
                 xleftmargin   = 0.8cm,
                 xrightmargin  = 0.8cm,
               ]{haskell}
    primes :: [Integer]
    primes = sieve [2..]
    
    sieve :: [Integer] -> [Integer]
    sieve (p:ns) = p : sieve [n | n <- ns, mod n p /= 0]
\end{minted}
\vspace*{-0.3cm}
\caption{Computing the prime numbers.}
\label{fig:primes.hs}
\end{figure}

Figure \ref{fig:primes.hs} on page \pageref{fig:primes.hs} shows a \textsl{Haskell} program to compute \textbf{all}
primes.  Yes, you have read that correct.  It doesn't compute the primes up to a given number, but rather it
computes \underline{all} primes.

The first thing to note is that line 1 and line 4 are type annotations.  They have only
been added to aid us in understanding the program.  If we would drop these lines, the program would still
work.  Hence, we have an efficient 2-line program to compute the prime numers.  Lets discuss this program line
by line. 
\begin{enumerate}[(a)]
\item Line 1 states that the function \texttt{primes} returns a list of \texttt{Integer}s.
  \texttt{Integer} is the type of all arbitrary precision integers.  The fact that we have enclosed
  the type name \texttt{Integer} in the square brackets ``\texttt{[}'' and ``\texttt{]}'' denotes that the
  result has the type \textsl{list} of \texttt{Integer}.
\item In line 2, the expression ``\texttt{[2..]}''  denotes the list of all integers starting from $2$.  Since
  \textsl{Haskell} is lazy, it is able to support infinite data structures.  The trick is that these lists are
  only evaluated as much as they are needed.  As long as we do not inspect the complete list, everything works
  fine. 
\item Line 2 calls the function \texttt{sieve} with the argument \texttt{[2..]}.  \textsl{Haskell} uses prefix
  notation for calling a function.  If $f$ is a function an $a_1$, $a_2$, and $a_3$ are arguments of this
  function, then the invocation of $f$ with these arguments is written as
  \\[0.2cm]
  \hspace*{1.3cm}
  $f$ $a_1$ $a_2$ $a_3$
  \\[0.2cm]
  \textbf{Note} that the expression
  \\[0.2cm]
  \hspace*{1.3cm}
  $f(a_1, a_2, a_3)$
  \\[0.2cm]
  denotes something different:  This expression would apply the function $f$ to a single argument, which is the
  triple $(a_1, a_2, a_3)$.
\item Line 4 declares the type of the function \texttt{sieve}.  This function takes one argument, which is a
  list of \texttt{Integer}s and returns a list of \texttt{Integer}s.
\item Line 5 defines the function \texttt{sieve} that takes a list of numbers $l$ that has the following
  properties:
  \begin{itemize}
  \item The list $l$ is a sorted ascendingly.
  \item If $p$ is the first element of the list $l$, then $p$ is a prime number.
  \item The list $l$ does not contain multiples of any number $q$ that is less than $p$:
    \\[0.2cm]
    \hspace*{1.3cm}
    $\forall q \in \mathbb{N}: \bigl(q < p \rightarrow \forall n \in \mathbb{N}: n \cdot q \not\in l\bigr)$.
  \end{itemize}
  Given a list $l$ with these properties, the expression
  \\[0.2cm]
  \hspace*{1.3cm}
  \texttt{sieve l}
  \\[0.2cm]
  returns a list of all prime number that are greater or equal than $p$, where $p$ is the first element of $l$: 
  \\[0.2cm]
  \hspace*{1.3cm}
  $\mathtt{sieve}\;l = [\; q \in l \mid q \in \mathbb{P}\;]$.
  \\[0.2cm]
  Hence, when \texttt{sieve} is called with the list of all natural number greater or equal than $2$
  it will return the list of all prime numbers, since $2$ is the smallest  prime numer.

  There is a lot going on in the definition of the function \texttt{sieve}.  We will discuss this definition
  in minute detail.
  \begin{enumerate}[1.]
  \item The function \texttt{sieve} is defined via \blue{matching}.  We will discuss matching in more detail in
    the next chapter.  For now we just mention that the expression
    \\[0.2cm]
    \hspace*{1.3cm}
    \texttt{(p:ns)}
    \\[0.2cm]
    matches a list with first element \texttt{p}.  The remaining elements are collected in the list
    \texttt{ns}.  For example, if \texttt{l = [2..]}, i.e.~if $l$ is the list of all natural numbers greater or
    equal than $2$, then the variable \texttt{p} is bound to the number $2$, while the variable \texttt{ns} is
    bound to the list of all natural numbers greater or equal than $3$, i.e.~\texttt{ns = [3..]}.
  \item The right hand side of the function definition, i.e.~the part after the symbol ``\texttt{=}'' defines the
    value that is computed by the function \texttt{sieve}.  This value is computed by calling the function
    ``\texttt{:}'', which is also known as the \blue{cons} function because it \underline{cons}tructs a list.
    The operator ``\texttt{:}'' takes two arguments.  The first argument is a value \texttt{u} of some type $a$
    and the second argument \texttt{us} is list of elements of the same type $a$.  An expression of the form
    \\[0.2cm]
    \hspace*{1.3cm}
    \texttt{u :\;us}
    \\[0.2cm]
    then returns a list where \texttt{u} is the first elements and \texttt{us} are the remaining elements.
    For example, we have
    \\[0.2cm]
    \hspace*{1.3cm}
    \texttt{1 :\;[2, 3, 4] = [1, 2, 3, 4]}.
  \item The recursive invocation of the function sieve takes a \blue{list comprehension} as ist first argument.
    the expression
    \\[0.2cm]
    \hspace*{1.3cm}
    \texttt{[n | n <- ns, mod n p /= 0]}
    \\[0.2cm]
    computes the list of all those number \texttt{n} from the list \texttt{ns} that have a non-zero remainder
    when divided by the prime number \texttt{p}, i.e.~this list contains all those number from the list
    \texttt{ns} that are not multiples of \texttt{p}.  There are two things to note here concerning the syntax
    of \textsl{Haskell}:
    \begin{itemize}
    \item Functions are written with prefix notation.  For example, we first write the function name
      \texttt{mode} followed by the arguments \texttt{m} and \texttt{p}.  In \textsl{Python} the expression
      \\[0.2cm]
      \hspace*{1.3cm}
      \texttt{mod m p}
      \\[0.2cm]
      would have been written as \ \texttt{m \% p}.
    \item The operator ``\texttt{/=}'' expresses inequality, i.e.~the \textsl{Haskell} expression
      \\[0.2cm]
      \hspace*{1.3cm}
      \texttt{a /= b} \qquad would be written as \qquad \texttt{a != b} 
      \\[0.2cm]
      in the programming language \textsl{Python}.
    \end{itemize}
    Putting everything together, the \textsl{Haskell} expression 
    \\[0.2cm]
    \hspace*{1.3cm}
    \texttt{[n | n <- ns, mod n p /= 0]}
    \\[0.2cm]
    is therefore equivalent to the \textsl{Python} expression
    \\[0.2cm]
    \hspace*{1.3cm}
    \texttt{[n for n in ns if n \% p != 0]}.
  \end{enumerate}
\end{enumerate}


\bibliographystyle{alpha}
\bibliography{cs}



\end{document}




%%% Local Variables:
%%% mode: latex
%%% TeX-master: t
%%% End:
